% القسم: المناهج
% الفصل: الاستقراء
% المبحث: الاستقراء البنيوي

\documentclass[../../../include/open-logic-section]{subfiles}

\begin{document}

\olfileid{mth}{ind}{sti}

\olsection{الاستقراء البنيوي}

جرى استعمال الاستقراء إلى الآن في أحكام تعم الأعداد الطبيعية. وله
نظير مباشر يثبت الأحكام لجميع !!{element} مجموعة معرّفة استقرائيًّا.
ويسمى هذا \emph{الاستقراء البنيوي}، لأن الحجة فيه تتبع بنية الأشياء
التي أنشأها التعريف الاستقرائي.

وقوام التعريف الاستقرائي أمران: (أ)~تعيين !!{element} «الأولية» في
المجموعة، و(ب)~تعيين العمليات التي تنشئ !!{element} جديدة منها.
فالأشياء الأولية في مثال الحدود السليمة هي الحروف، والعملية الوحيدة
هي:
\begin{align*}
  {\OLId{latin-ordinary}{o}}({\OLId{latin-ordinary}{s}}_{\OLMathNumeral{1}}, {\OLId{latin-ordinary}{s}}_{\OLMathNumeral{2}}) &= [{\OLId{latin-ordinary}{s}}_{\OLMathNumeral{1}} \circ {\OLId{latin-ordinary}{s}}_{\OLMathNumeral{2}}]
\end{align*}
ويجوز أن تُرى الأعداد الطبيعية~$\Nat$ نفسها ثمرة تعريف استقرائي:
فالابتداء بـ~${\OLMathNumeral{0}}$، والعملية دالة الخلف~${\OLId{latin-ordinary}{x}} + {\OLMathNumeral{1}}$.

وإذا أُريد إثبات حكم في جميع أفراد مجموعة معرّفة استقرائيًّا، أي
إثبات أن كل !!{element} فيها ذو خاصية~${\OLId{latin-looped}{P}}$، وجب أمران:
\begin{enumerate}
\item إثبات الخاصية~${\OLId{latin-looped}{P}}$ للأشياء الأولية.
\item وإثبات، لكل عملية~${\OLId{latin-ordinary}{o}}$، أن اتصاف معاملاتها بـ~${\OLId{latin-looped}{P}}$ يوجب اتصاف
  نتيجتها بها أيضًا.
\end{enumerate}
فالحكم على جميع الحدود السليمة، مثلًا، يثبت أولًا لكل حرف، ثم يثبت
لـ$[{\OLId{latin-ordinary}{s}}_{\OLMathNumeral{1}} \circ {\OLId{latin-ordinary}{s}}_{\OLMathNumeral{2}}]$ متى ثبت لكل من ${\OLId{latin-ordinary}{s}}_{\OLMathNumeral{1}}$ و${\OLId{latin-ordinary}{s}}_{\OLMathNumeral{2}}$ منفردًا.

\begin{prop}
  في كل حد سليم~${\OLId{latin-ordinary}{t}}$، يساوي عدد الرموز $[$ عدد الرموز $]$.
\end{prop}

\begin{proof}
نستعمل الاستقراء البنيوي؛ فالحدود السليمة معرّفة استقرائيًّا، والحروف
أفرادها الأولية، و${\OLId{latin-ordinary}{o}}$ هي العملية المنشئة للحد الجديد من حدين سابقين.
\begin{enumerate}
\item القضية صادقة لكل حرف؛ فعدد الرموز $[$ في الحرف المنفرد~${\OLMathNumeral{0}}$،
  وعدد الرموز $]$ فيه~${\OLMathNumeral{0}}$ أيضًا.
\item افرض أن عدد الرموز $[$ في ${\OLId{latin-ordinary}{s}}_{\OLMathNumeral{1}}$ يساوي عدد الرموز $]$ في
  ${\OLId{latin-ordinary}{s}}_{\OLMathNumeral{1}}$، وأن النظير صادق لـ${\OLId{latin-ordinary}{s}}_{\OLMathNumeral{2}}$. فعدد الرموز $[$ في
  ${\OLId{latin-ordinary}{o}}({\OLId{latin-ordinary}{s}}_{\OLMathNumeral{1}}, {\OLId{latin-ordinary}{s}}_{\OLMathNumeral{2}})$، أي في $[{\OLId{latin-ordinary}{s}}_{\OLMathNumeral{1}} \circ {\OLId{latin-ordinary}{s}}_{\OLMathNumeral{2}}]$، هو مجموع عدديها في
  ${\OLId{latin-ordinary}{s}}_{\OLMathNumeral{1}}$ و${\OLId{latin-ordinary}{s}}_{\OLMathNumeral{2}}$ مع زيادة واحد. وعدد الرموز $]$ في
  ${\OLId{latin-ordinary}{o}}({\OLId{latin-ordinary}{s}}_{\OLMathNumeral{1}}, {\OLId{latin-ordinary}{s}}_{\OLMathNumeral{2}})$ هو كذلك مجموع عدديها في ${\OLId{latin-ordinary}{s}}_{\OLMathNumeral{1}}$ و${\OLId{latin-ordinary}{s}}_{\OLMathNumeral{2}}$ مع زيادة
  واحد. ومن ثم يتساوى عدد الرموز $[$ في ${\OLId{latin-ordinary}{o}}({\OLId{latin-ordinary}{s}}_{\OLMathNumeral{1}}, {\OLId{latin-ordinary}{s}}_{\OLMathNumeral{2}})$ وعدد الرموز
  $]$ فيه.
\end{enumerate}
\end{proof}

\begin{prob}
  أثبت بالاستقراء البنيوي أنه لا يبدأ أي حد سليم بالرمز~$]$.
\end{prob}

ولنأت ببرهان بنيوي آخر. القطعة الابتدائية الصحيحة من سلسلة الرموز~${\OLId{latin-ordinary}{t}}$
هي سلسلة~${\OLId{latin-ordinary}{s}}$ تطابق ${\OLId{latin-ordinary}{t}}$ رمزًا رمزًا من جهة اليسار، على أن يكون ${\OLId{latin-ordinary}{t}}$
أطول منها. فـ$[{\OLId{latin-ordinary}{a}} \circ {}$، مثلًا، قطعة ابتدائية صحيحة من
$[{\OLId{latin-ordinary}{a}} \circ {\OLId{latin-ordinary}{b}}]$؛ وليس $[{\OLId{latin-ordinary}{b}} \circ {}$ كذلك لاختلاف الرمز الثاني، ولا
$[{\OLId{latin-ordinary}{a}} \circ {\OLId{latin-ordinary}{b}}]$ نفسها لتساوي الطولين.

\begin{prop}\ollabel{prop:initial}
  في كل قطعة ابتدائية صحيحة غير خالية من حد سليم~${\OLId{latin-ordinary}{t}}$ عدد من الرموز $[$ أكبر من
  عدد الرموز $]$.
\end{prop}

\begin{proof}
  نستعمل الاستقراء على~${\OLId{latin-ordinary}{t}}$:
  \begin{enumerate}
  \item ${\OLId{latin-ordinary}{t}}$ حرف بمفرده: عندئذ ليست لـ${\OLId{latin-ordinary}{t}}$ أي قطعة ابتدائية صحيحة.
  \item ${\OLId{latin-ordinary}{t}} = [{\OLId{latin-ordinary}{s}}_{\OLMathNumeral{1}} \circ {\OLId{latin-ordinary}{s}}_{\OLMathNumeral{2}}]$ لبعض حدين سليمين ${\OLId{latin-ordinary}{s}}_{\OLMathNumeral{1}}$ و~${\OLId{latin-ordinary}{s}}_{\OLMathNumeral{2}}$. إذا
    كان ${\OLId{latin-ordinary}{r}}$ قطعة ابتدائية صحيحة من ${\OLId{latin-ordinary}{t}}$، فثمة عدة احتمالات:
    \begin{enumerate}
    \item ${\OLId{latin-ordinary}{r}}$ هي $[$ فقط: عندئذ في ${\OLId{latin-ordinary}{r}}$ رمز $[$ واحد أكثر من الرموز
      $]$.
    \item ${\OLId{latin-ordinary}{r}}$ هي $[{\OLId{latin-ordinary}{r}}_{\OLMathNumeral{1}}$، حيث ${\OLId{latin-ordinary}{r}}_{\OLMathNumeral{1}}$ قطعة ابتدائية صحيحة من~${\OLId{latin-ordinary}{s}}_{\OLMathNumeral{1}}$:
      بما أن ${\OLId{latin-ordinary}{s}}_{\OLMathNumeral{1}}$ حد سليم، فإن ${\OLId{latin-ordinary}{r}}_{\OLMathNumeral{1}}$، بحسب فرض الاستقراء، تحتوي على
      رموز $[$ أكثر من الرموز $]$؛ وكذلك $[{\OLId{latin-ordinary}{r}}_{\OLMathNumeral{1}}$.
    \item ${\OLId{latin-ordinary}{r}}$ هي $[{\OLId{latin-ordinary}{s}}_{\OLMathNumeral{1}}$ أو $[{\OLId{latin-ordinary}{s}}_{\OLMathNumeral{1}} \circ {}$: بحسب النتيجة السابقة،
      يتساوى عدد الرموز $[$ و$]$ في~${\OLId{latin-ordinary}{s}}_{\OLMathNumeral{1}}$؛ ولذلك يزيد عدد الرموز $[$
      في $[{\OLId{latin-ordinary}{s}}_{\OLMathNumeral{1}}$ أو $[{\OLId{latin-ordinary}{s}}_{\OLMathNumeral{1}} \circ {}$ بواحد على عدد الرموز~$]$.
    \item ${\OLId{latin-ordinary}{r}}$ هي $[{\OLId{latin-ordinary}{s}}_{\OLMathNumeral{1}} \circ {\OLId{latin-ordinary}{r}}_{\OLMathNumeral{2}}$، حيث ${\OLId{latin-ordinary}{r}}_{\OLMathNumeral{2}}$ قطعة ابتدائية صحيحة
      من~${\OLId{latin-ordinary}{s}}_{\OLMathNumeral{2}}$: بحسب فرض الاستقراء، تحتوي ${\OLId{latin-ordinary}{r}}_{\OLMathNumeral{2}}$ على رموز $[$ أكثر من
      الرموز $]$. وبحسب النتيجة السابقة، يتساوى عدد الرموز $[$ و$]$
      في~${\OLId{latin-ordinary}{s}}_{\OLMathNumeral{1}}$. لذلك يكون عدد الرموز $[$ في $[{\OLId{latin-ordinary}{s}}_{\OLMathNumeral{1}} \circ {\OLId{latin-ordinary}{r}}_{\OLMathNumeral{2}}$ أكبر من
      عدد الرموز~$]$.
    \item ${\OLId{latin-ordinary}{r}}$ هي $[{\OLId{latin-ordinary}{s}}_{\OLMathNumeral{1}} \circ {\OLId{latin-ordinary}{s}}_{\OLMathNumeral{2}}$: بحسب النتيجة السابقة، يتساوى عدد
      الرموز $[$ و$]$ في ${\OLId{latin-ordinary}{s}}_{\OLMathNumeral{1}}$، ويصدق الشيء نفسه على~${\OLId{latin-ordinary}{s}}_{\OLMathNumeral{2}}$. ولذلك
      يوجد في $[{\OLId{latin-ordinary}{s}}_{\OLMathNumeral{1}} \circ {\OLId{latin-ordinary}{s}}_{\OLMathNumeral{2}}$ رمز $[$ واحد أكثر من الرموز~$]$.
    \end{enumerate}
  \end{enumerate}
\end{proof}

\end{document}
